%% Écrit par outils/lexique.py — ne pas modifier à la main
\begin{lexique}
\lexlettre{A}
\lexmot{across from}{əˈkrɑːs frəm}{en face de \registre{américain}}{9}
\lexmot{always}{ˈɑːlweɪz}{toujours}{6}
\lexmot{And you?}{ənd ˈjuː}{Et toi ? / Et vous ?}{1}
\lexmot{Are you free on Friday?}{ɑːr juː ˈfriː ɑːn ˈfraɪdeɪ}{Tu es libre vendredi ?}{15}
\lexmot{Are you French?}{ɑːr juː ˈfrentʃ}{Tu es français(e) ? / Vous êtes français(e) ?}{2}
\lexmot{Are you ready to order?}{ɑːr juː ˌredi tu ˈɔːrdɚ}{Vous avez choisi ?}{11}
\lexmot{arrive}{əˈraɪv}{arriver}{13}
\lexmot{at}{æt}{arobase (@)}{3}
\lexmot{at six}{æt ˈsɪks}{à six heures}{6}
\lexmot{at the moment}{ət ðə ˈmoʊmənt}{en ce moment}{13}
\lexmot{aunt · uncle}{ænt · ˈʌŋkəl}{la tante · l'oncle}{4}
\lexmot{Australia · Australian}{ɑːˈstreɪljə · ɑːˈstreɪljən}{l'Australie · australien(ne)}{2}
\lexmot{autumn}{ˈɔːtəm}{l'automne \registre{britannique}}{14}
\lexmot{an avenue}{ən ˈævənuː}{une avenue}{9}
\lexlettre{B}
\lexmot{a baby}{ə ˈbeɪbi}{un bébé}{4}
\lexmot{a bag}{ə ˈbæɡ}{un sac}{5}
\lexmot{a balcony}{ə ˈbælkəni}{un balcon}{8}
\lexmot{a band}{ə ˈbænd}{un groupe (de musique)}{12}
\lexmot{a bank}{ə ˈbæŋk}{une banque}{9}
\lexmot{a bathroom}{ə ˈbæθruːm}{une salle de bains}{8}
\lexmot{the beach}{ðə ˈbiːtʃ}{la plage}{16}
\lexmot{a beautiful city}{ə ˌbjuːtɪfəl ˈsɪti}{une belle ville}{2}
\lexmot{a bed}{ə ˈbed}{un lit}{8}
\lexmot{a bedroom}{ə ˈbedruːm}{une chambre}{8}
\lexmot{a beginner}{ə bɪˈɡɪnɚ}{un débutant, une débutante}{2}
\lexmot{Best wishes!}{ˌbest ˈwɪʃɪz}{Tous mes vœux !}{20}
\lexmot{a big city}{ə ˌbɪɡ ˈsɪti}{une grande ville}{2}
\lexmot{a birthday party}{ə ˈbɝːθdeɪ ˌpɑːrti}{une fête d'anniversaire}{20}
\lexmot{black}{blæk}{noir \otNuance{1C1C1C}}{5}
\lexmot{blow out the candles}{ˌbloʊ ˌaʊt ðə ˈkændəlz}{souffler les bougies}{20}
\lexmot{blue}{bluː}{bleu \otNuance{2F6FD6}}{5}
\lexmot{a book}{ə ˈbʊk}{un livre}{5}
\lexmot{boots}{buːts}{des bottes}{13}
\lexmot{a bottle}{ə ˈbɑːtəl}{une bouteille}{5}
\lexmot{a box}{ə ˈbɑːks}{une boîte, un carton}{5}
\lexmot{a boy · a girl}{ə ˈbɔɪ · ə ˈɡɝːl}{un garçon · une fille, une jeune femme}{4}
\lexmot{Brazil · Brazilian}{brəˈzɪl · brəˈzɪliən}{le Brésil · brésilien(ne)}{2}
\lexmot{bread}{bred}{le pain}{7}
\lexmot{a bridge}{ə ˈbrɪdʒ}{un pont}{9}
\lexmot{bring}{brɪŋ}{apporter, amener}{15}
\lexmot{brown}{braʊn}{marron \otNuance{8A5A3B}}{5}
\lexmot{brunch}{brʌntʃ}{le brunch}{18}
\lexmot{a building}{ə ˈbɪldɪŋ}{un immeuble, un bâtiment}{8}
\lexmot{a burger}{ə ˈbɝːɡɚ}{un hamburger}{11}
\lexmot{a bus stop}{ə ˈbʌs ˌstɑːp}{un arrêt de bus}{9}
\lexmot{busy}{ˈbɪzi}{occupé}{15}
\lexmot{buy · bought}{baɪ · bɑːt}{acheter}{19}
\lexmot{bye}{baɪ}{salut, au revoir}{1}
\lexlettre{C}
\lexmot{a cake}{ə ˈkeɪk}{un gâteau}{20}
\lexmot{call}{kɑːl}{appeler, téléphoner à}{13}
\lexmot{call · called}{kɑːl · kɑːld}{appeler}{18}
\lexmot{call back}{ˌkɑːl ˈbæk}{rappeler}{13}
\lexmot{Can I have a coffee, please?}{kən aɪ ˌhæv ə ˈkɑːfi pliːz}{Je peux avoir un café, s'il vous plaît ?}{11}
\lexmot{Can I help you?}{kən aɪ ˈhelp juː}{Je peux vous aider ?}{10}
\lexmot{Can I leave a message?}{kən aɪ ˌliːv ə ˈmesɪdʒ}{Puis-je laisser un message ?}{13}
\lexmot{Can I pay by card?}{kən aɪ ˌpeɪ baɪ ˈkɑːrd}{Je peux payer par carte ?}{10}
\lexmot{Can I speak to…?}{kən aɪ ˈspiːk tə}{Pourrais-je parler à… ?}{13}
\lexmot{Can I take a message?}{kən aɪ ˌteɪk ə ˈmesɪdʒ}{Puis-je prendre un message ?}{13}
\lexmot{Can I try it on?}{kən aɪ ˌtraɪ ɪt ˈɑːn}{Je peux l'essayer ?}{10}
\lexmot{Can we have the check, please?}{kən wiː ˌhæv ðə ˈtʃek pliːz}{L'addition, s'il vous plaît.}{11}
\lexmot{Can you swim?}{kən juː ˈswɪm}{Tu sais nager ?}{12}
\lexmot{Canada · Canadian}{ˈkænədə · kəˈneɪdiən}{le Canada · canadien(ne)}{2}
\lexmot{a candle}{ə ˈkændəl}{une bougie}{20}
\lexmot{a cap}{ə ˈkæp}{une casquette}{13}
\lexmot{a card}{ə ˈkɑːrd}{une carte}{20}
\lexmot{carry}{ˈkæri}{porter (dans ses bras)}{13}
\lexmot{cash}{kæʃ}{en espèces, du liquide}{10}
\lexmot{a cat}{ə ˈkæt}{un chat}{4}
\lexmot{celebrate}{ˈseləbreɪt}{fêter, célébrer}{20}
\lexmot{a cent}{ə ˈsent}{un cent (centième de dollar)}{10}
\lexmot{a century}{ə ˈsentʃəri}{un siècle}{17}
\lexmot{a chair}{ə ˈtʃer}{une chaise}{5}
\lexmot{cheap · expensive}{tʃiːp · ɪkˈspensɪv}{bon marché · cher}{10}
\lexmot{the check}{ðə ˈtʃek}{l'addition \registre{américain}}{11}
\lexmot{check in}{ˌtʃek ˈɪn}{s'enregistrer (à l'hôtel)}{19}
\lexmot{Cheers!}{tʃɪrz}{Santé ! (pour trinquer)}{20}
\lexmot{cheese}{tʃiːz}{le fromage}{7}
\lexmot{chicken}{ˈtʃɪkɪn}{le poulet}{7}
\lexmot{children}{ˈtʃɪldrən}{les enfants (singulier : \en{child})}{4}
\lexmot{China · Chinese}{ˈtʃaɪnə · tʃaɪˈniːz}{la Chine · chinois(e)}{2}
\lexmot{chocolate}{ˈtʃɑːklət}{le chocolat}{7}
\lexmot{clean · cleaned}{kliːn · kliːnd}{nettoyer}{18}
\lexmot{a closet}{ə ˈklɑːzɪt}{un placard, une penderie}{8}
\lexmot{clothes}{kloʊðz}{les vêtements}{10}
\lexmot{a coat}{ə ˈkoʊt}{un manteau}{10}
\lexmot{a coffee}{ə ˈkɑːfi}{un café}{11}
\lexmot{coffee}{ˈkɑːfi}{le café}{7}
\lexmot{come · came}{kʌm · keɪm}{venir}{19}
\lexmot{Congratulations!}{kənˌɡrætʃəˈleɪʃənz}{Félicitations !}{20}
\lexmot{cook · cooked}{kʊk · kʊkt}{cuisiner}{18}
\lexmot{cook dinner}{ˌkʊk ˈdɪnɚ}{préparer le dîner}{6}
\lexmot{cooking}{ˈkʊkɪŋ}{la cuisine, cuisiner}{7}
\lexmot{Could I have the fish?}{kʊd aɪ ˌhæv ðə ˈfɪʃ}{Pourrais-je avoir le poisson ?}{11}
\lexmot{cousin}{ˈkʌzən}{le cousin, la cousine}{4}
\lexmot{cross the street}{ˌkrɑːs ðə ˈstriːt}{traverser la rue}{9}
\lexmot{crowded}{ˈkraʊdɪd}{bondé, plein de monde}{17}
\lexmot{a cup}{ə ˈkʌp}{une tasse}{5}
\lexlettre{D}
\lexmot{dance}{dæns}{danser}{12}
\lexmot{dance · danced}{dæns · dænst}{danser}{18}
\lexmot{dancing}{ˈdænsɪŋ}{la danse, danser}{7}
\lexmot{daughter · son}{ˈdɑːtɚ · sʌn}{la fille (l'enfant) · le fils}{4}
\lexmot{degrees}{dɪˈɡriːz}{degrés}{14}
\lexmot{a desk}{ə ˈdesk}{un bureau (le meuble)}{5}
\lexmot{dessert}{dɪˈzɝːt}{le dessert}{11}
\lexmot{Did you have a nice weekend?}{dɪd juː ˌhæv ə ˌnaɪs ˈwiːkend}{Tu as passé un bon week-end ?}{18}
\lexmot{Do you want to come?}{də jə ˈwɑːnt tə ˈkʌm}{Tu veux venir ?}{7}
\lexmot{a dog}{ə ˈdɑːɡ}{un chien}{4}
\lexmot{Don't worry.}{ˌdoʊnt ˈwɝːi}{Ne t'inquiète pas. / Ne vous inquiétez pas.}{2}
\lexmot{a door}{ə ˈdɔːr}{une porte}{8}
\lexmot{dot}{dɑːt}{point (dans une adresse)}{3}
\lexmot{draw}{drɑː}{dessiner}{12}
\lexmot{a dress}{ə ˈdres}{une robe}{10}
\lexmot{drink a coffee}{ˌdrɪŋk ə ˈkɑːfi}{boire un café}{6}
\lexmot{drive a car}{ˌdraɪv ə ˈkɑːr}{conduire une voiture}{12}
\lexlettre{E}
\lexmot{early · late}{ˈɝːli · leɪt}{tôt, en avance · tard, en retard}{6}
\lexmot{eat · ate}{iːt · eɪt}{manger}{19}
\lexmot{eggs}{eɡz}{les œufs}{7}
\lexmot{an elevator}{ən ˈeləveɪtɚ}{un ascenseur}{8}
\lexmot{an email address}{ən ˈiːmeɪl ˌædres}{une adresse e-mail}{3}
\lexmot{England · English}{ˈɪŋɡlənd · ˈɪŋɡlɪʃ}{l'Angleterre · anglais(e)}{2}
\lexmot{every day}{ˌevri ˈdeɪ}{tous les jours}{6}
\lexmot{Excuse me.}{ɪkˈskjuːz miː}{Excusez-moi.}{9}
\lexlettre{F}
\lexmot{fall}{fɑːl}{l'automne \registre{américain}}{14}
\lexmot{family}{ˈfæməli}{la famille}{4}
\lexmot{famous}{ˈfeɪməs}{célèbre}{17}
\lexmot{fast}{fæst}{vite}{12}
\lexmot{finish work}{ˌfɪnɪʃ ˈwɝːk}{finir le travail}{6}
\lexmot{the first floor}{ðə ˌfɝːst ˈflɔːr}{le rez-de-chaussée (américain)}{8}
\lexmot{first name}{ˈfɝːst ˌneɪm}{prénom}{3}
\lexmot{fish}{fɪʃ}{le poisson}{7}
\lexmot{a flight}{ə ˈflaɪt}{un vol}{16}
\lexmot{fly}{flaɪ}{prendre l'avion, voler}{16}
\lexmot{For here or to go?}{fɚ ˈhɪr ɔːr tə ˈɡoʊ}{Sur place ou à emporter ?}{11}
\lexmot{France · French}{fræns · frentʃ}{la France · français(e)}{2}
\lexmot{free}{friː}{libre ; gratuit}{15}
\lexmot{French food}{ˌfrentʃ ˈfuːd}{la cuisine française}{7}
\lexmot{Friday}{ˈfraɪdeɪ}{vendredi}{6}
\lexmot{a fridge}{ə ˈfrɪdʒ}{un réfrigérateur}{8}
\lexmot{fries}{fraɪz}{des frites}{11}
\lexmot{fruit}{fruːt}{les fruits}{7}
\lexlettre{G}
\lexmot{a game}{ə ˈɡeɪm}{un match, un jeu}{12}
\lexmot{a garden}{ə ˈɡɑːrdən}{un jardin}{8}
\lexmot{Germany · German}{ˈdʒɝːməni · ˈdʒɝːmən}{l'Allemagne · allemand(e)}{2}
\lexmot{get · got}{ɡet · ɡɑːt}{obtenir ; arriver}{19}
\lexmot{get off}{ˌɡet ˈɑːf}{descendre (du métro, du bus)}{9}
\lexmot{get up}{ˌɡet ˈʌp}{se lever}{6}
\lexmot{a gift}{ə ˈɡɪft}{un cadeau}{17}
\lexmot{give · gave}{ɡɪv · ɡeɪv}{donner}{19}
\lexmot{a glass of water}{ə ˌɡlæs əv ˈwɑːtɚ}{un verre d'eau}{11}
\lexmot{glasses}{ˈɡlæsɪz}{des lunettes}{5}
\lexmot{gloves}{ɡlʌvz}{des gants}{10}
\lexmot{go · went}{ɡoʊ · went}{aller}{19}
\lexmot{go camping}{ˌɡoʊ ˈkæmpɪŋ}{faire du camping}{16}
\lexmot{go home}{ˌɡoʊ ˈhoʊm}{rentrer à la maison}{6}
\lexmot{go out}{ˌɡoʊ ˈaʊt}{sortir}{15}
\lexmot{go straight ahead}{ˌɡoʊ ˌstreɪt əˈhed}{aller tout droit}{9}
\lexmot{go to bed}{ˌɡoʊ tə ˈbed}{aller se coucher}{6}
\lexmot{go to the movies}{ˌɡoʊ tə ðə ˈmuːviz}{aller au cinéma}{15}
\lexmot{go to work}{ˌɡoʊ tə ˈwɝːk}{aller au travail}{6}
\lexmot{good afternoon}{ɡʊd ˌæftɚˈnuːn}{bonjour (l'après-midi)}{1}
\lexmot{good evening}{ɡʊd ˈiːvnɪŋ}{bonsoir (en arrivant)}{1}
\lexmot{good morning}{ɡʊd ˈmɔːrnɪŋ}{bonjour (le matin)}{1}
\lexmot{good night}{ɡʊd ˈnaɪt}{bonne nuit, bonne soirée (en partant tard, ou avant de dormir)}{1}
\lexmot{goodbye}{ɡʊdˈbaɪ}{au revoir}{1}
\lexmot{Got it!}{ˈɡɑːt ɪt}{C'est noté ! Compris !}{3}
\lexmot{grandmother · grandfather}{ˈɡrænˌmʌðɚ · ˈɡrænˌfɑːðɚ}{la grand-mère · le grand-père}{4}
\lexmot{grandparents}{ˈɡrænˌperənts}{les grands-parents}{4}
\lexmot{gray}{ɡreɪ}{gris \otNuance{9AA0A6} (britannique : \en{grey})}{5}
\lexmot{Great idea!}{ˌɡreɪt aɪˈdiːə}{Bonne idée !}{15}
\lexmot{green}{ɡriːn}{vert \otNuance{3B9B4E}}{5}
\lexlettre{H}
\lexmot{Happy birthday!}{ˌhæpi ˈbɝːθdeɪ}{Joyeux anniversaire !}{20}
\lexmot{the harbor}{ðə ˈhɑːrbɚ}{le port}{19}
\lexmot{a hat}{ə ˈhæt}{un chapeau, un bonnet}{10}
\lexmot{hate}{heɪt}{détester}{7}
\lexmot{have · had}{hæv · hæd}{avoir}{19}
\lexmot{have a nice day}{ˌhæv ə ˌnaɪs ˈdeɪ}{bonne journée}{1}
\lexmot{have a nice evening}{ˌhæv ə ˌnaɪs ˈiːvnɪŋ}{bonne soirée}{1}
\lexmot{have breakfast}{ˌhæv ˈbrekfəst}{prendre le petit-déjeuner}{6}
\lexmot{have lunch}{ˌhæv ˈlʌntʃ}{déjeuner}{6}
\lexmot{hello}{həˈloʊ}{bonjour (à toute heure)}{1}
\lexmot{Here you are.}{ˌhɪr juː ˈɑːr}{Voici. Tenez.}{11}
\lexmot{He's in a meeting.}{hiz ɪn ə ˈmiːtɪŋ}{Il est en réunion.}{13}
\lexmot{He's so cute!}{hiːz ˌsoʊ ˈkjuːt}{Il est trop mignon !}{4}
\lexmot{hey}{heɪ}{salut \registre{familier}}{1}
\lexmot{hi}{haɪ}{salut, bonjour}{1}
\lexmot{history}{ˈhɪstəri}{l'histoire (le passé)}{17}
\lexmot{Hold on, please.}{ˌhoʊld ˈɑːn ˌpliːz}{Ne quittez pas.}{13}
\lexmot{a holiday}{ə ˈhɑːlədeɪ}{un jour férié, une fête}{14}
\lexmot{a hoodie}{ə ˈhʊdi}{un sweat à capuche}{13}
\lexmot{a hospital}{ə ˈhɑːspɪtəl}{un hôpital}{9}
\lexmot{a hotel}{ə hoʊˈtel}{un hôtel}{9}
\lexmot{a hotel room}{ə hoʊˈtel ˌruːm}{une chambre d'hôtel}{19}
\lexmot{a house}{ə ˈhaʊs}{une maison}{8}
\lexmot{How are you?}{ˌhaʊ ˈɑːr juː}{Comment ça va ? / Comment allez-vous ?}{1}
\lexmot{How do I get to the museum?}{ˌhaʊ duː aɪ ˈɡet tə ðə mjuːˈziːəm}{Comment je vais au musée ?}{9}
\lexmot{How do you spell it?}{haʊ duː juː ˈspel ɪt}{Comment ça s'écrit ?}{3}
\lexmot{How much are they?}{ˌhaʊ ˈmʌtʃ ɑːr ðeɪ}{Combien coûtent-ils ?}{10}
\lexmot{How much is it?}{ˌhaʊ ˈmʌtʃ ɪz ɪt}{Combien ça coûte ?}{10}
\lexmot{How old are you?}{ˌhaʊ ˈoʊld ɑːr juː}{Quel âge as-tu ? / Quel âge avez-vous ?}{3}
\lexmot{How was it?}{ˌhaʊ ˈwʌz ɪt}{C'était comment ?}{17}
\lexmot{How was your trip?}{ˌhaʊ wəz jʊr ˈtrɪp}{C'était comment, ton voyage ?}{19}
\lexmot{hurry up}{ˌhɝːi ˈʌp}{se dépêcher}{13}
\lexlettre{I}
\lexmot{I can sing.}{aɪ kən ˈsɪŋ}{Je sais chanter.}{12}
\lexmot{I can't ski.}{aɪ ˈkænt ˈskiː}{Je ne sais pas skier.}{12}
\lexmot{I didn't do much.}{aɪ ˌdɪdənt duː ˈmʌtʃ}{Je n'ai pas fait grand-chose.}{18}
\lexmot{I don't like it.}{aɪ ˌdoʊnt ˈlaɪk ɪt}{Je n'aime pas ça.}{7}
\lexmot{I live in an apartment.}{aɪ ˌlɪv ɪn ən əˈpɑːrtmənt}{J'habite dans un appartement.}{8}
\lexmot{I speak English.}{aɪ ˌspiːk ˈɪŋɡlɪʃ}{Je parle anglais.}{2}
\lexmot{I was at home.}{aɪ wəz ət ˈhoʊm}{J'étais à la maison.}{17}
\lexmot{I was born in France.}{aɪ wəz ˌbɔːrn ɪn ˈfræns}{Je suis né(e) en France.}{17}
\lexmot{ice cream}{ˈaɪs ˌkriːm}{de la glace}{11}
\lexmot{I'd like a salad.}{aɪd ˌlaɪk ə ˈsæləd}{Je voudrais une salade.}{11}
\lexmot{I'd love to!}{aɪd ˈlʌv tuː}{Avec plaisir !}{7}
\lexmot{I'll have the burger.}{aɪl ˌhæv ðə ˈbɝːɡɚ}{Je prendrai le burger.}{11}
\lexmot{I'll take it.}{aɪl ˈteɪk ɪt}{Je le prends.}{10}
\lexmot{I'm busy.}{aɪm ˈbɪzi}{Je suis occupé.}{15}
\lexmot{I'm fine, thanks.}{aɪm ˈfaɪn θæŋks}{Ça va bien, merci.}{1}
\lexmot{I'm from France.}{aɪm frəm ˈfræns}{Je viens de France.}{2}
\lexmot{I'm going to travel.}{aɪm ˌɡoʊɪŋ tə ˈtrævəl}{Je vais voyager.}{16}
\lexmot{I'm good, thanks.}{aɪm ˈɡʊd θæŋks}{Ça va bien, merci. \registre{très courant aux États-Unis}}{1}
\lexmot{I'm Léa.}{aɪm ˈleɪə}{Moi, c'est Léa. / Je suis Léa.}{1}
\lexmot{I'm not going to work.}{aɪm ˌnɑːt ˌɡoʊɪŋ tə ˈwɝːk}{Je ne vais pas travailler.}{16}
\lexmot{I'm OK.}{aɪm ˌoʊˈkeɪ}{Ça va.}{1}
\lexmot{I'm proud of you!}{aɪm ˈpraʊd əv juː}{Je suis fier de toi !}{20}
\lexmot{I'm sorry, I can't.}{aɪm ˈsɑːri aɪ ˈkænt}{Désolé, je ne peux pas.}{15}
\lexmot{I'm twenty-seven.}{aɪm ˌtwentiˈsevən}{J'ai vingt-sept ans.}{3}
\lexmot{I'm very well, thank you.}{aɪm ˌveri ˈwel ˈθæŋk juː}{Je vais très bien, merci.}{1}
\lexmot{an immigrant}{ən ˈɪmɪɡrənt}{un immigrant}{17}
\lexmot{improve}{ɪmˈpruːv}{progresser, s'améliorer}{20}
\lexmot{in 1990}{ɪn ˌnaɪntiːn ˈnaɪnti}{en 1990}{17}
\lexmot{in the afternoon}{ɪn ði ˌæftɚˈnuːn}{l'après-midi}{6}
\lexmot{in the evening}{ɪn ði ˈiːvnɪŋ}{le soir}{6}
\lexmot{in the morning}{ɪn ðə ˈmɔːrnɪŋ}{le matin}{6}
\lexmot{in two weeks}{ɪn ˌtuː ˈwiːks}{dans deux semaines}{16}
\lexmot{invite}{ɪnˈvaɪt}{inviter}{15}
\lexmot{Ireland · Irish}{ˈaɪrlənd · ˈaɪrɪʃ}{l'Irlande · irlandais(e)}{2}
\lexmot{Is it far?}{ɪz ɪt ˈfɑːr}{C'est loin ?}{9}
\lexmot{Is the tip included?}{ɪz ðə ˌtɪp ɪnˈkluːdɪd}{Le pourboire est compris ?}{11}
\lexmot{It was amazing!}{ɪt wəz əˈmeɪzɪŋ}{C'était génial !}{19}
\lexmot{It was fun!}{ɪt wəz ˈfʌn}{C'était sympa !}{18}
\lexmot{Italy · Italian}{ˈɪtəli · ɪˈtæljən}{l'Italie · italien(ne)}{2}
\lexmot{It's a five-minute walk.}{ɪts ə ˌfaɪv ˌmɪnɪt ˈwɑːk}{C'est à cinq minutes à pied.}{9}
\lexmot{It's cloudy.}{ɪts ˈklaʊdi}{Le ciel est nuageux.}{14}
\lexmot{It's cold.}{ɪts ˈkoʊld}{Il fait froid.}{14}
\lexmot{It's cool.}{ɪts ˈkuːl}{Il fait frais.}{14}
\lexmot{It's delicious!}{ɪts dɪˈlɪʃəs}{C'est délicieux !}{11}
\lexmot{It's foggy.}{ɪts ˈfɑːɡi}{Il y a du brouillard.}{14}
\lexmot{It's freezing!}{ɪts ˈfriːzɪŋ}{Il gèle ! Il fait un froid glacial !}{14}
\lexmot{It's hot.}{ɪts ˈhɑːt}{Il fait chaud.}{14}
\lexmot{It's near here.}{ɪts ˌnɪr ˈhɪr}{C'est près d'ici.}{9}
\lexmot{It's nine fifteen.}{ɪts ˌnaɪn fɪfˈtiːn}{Il est neuf heures et quart.}{6}
\lexmot{It's nine o'clock.}{ɪts ˌnaɪn əˈklɑːk}{Il est neuf heures.}{6}
\lexmot{It's nine thirty.}{ɪts ˌnaɪn ˈθɝːti}{Il est neuf heures et demie.}{6}
\lexmot{It's on March 15th.}{ɪts ɑːn ˌmɑːrtʃ fɪfˈtiːnθ}{C'est le 15 mars.}{14}
\lexmot{It's raining.}{ɪts ˈreɪnɪŋ}{Il pleut.}{14}
\lexmot{It's snowing.}{ɪts ˈsnoʊɪŋ}{Il neige.}{14}
\lexmot{It's sunny.}{ɪts ˈsʌni}{Il fait soleil.}{14}
\lexmot{It's twenty dollars.}{ɪts ˌtwenti ˈdɑːlɚz}{Ça coûte vingt dollars.}{10}
\lexmot{It's warm.}{ɪts ˈwɔːrm}{Il fait bon, il fait doux.}{14}
\lexmot{It's windy.}{ɪts ˈwɪndi}{Il y a du vent.}{14}
\lexlettre{J}
\lexmot{a jacket}{ə ˈdʒækɪt}{une veste, un blouson}{10}
\lexmot{Japan · Japanese}{dʒəˈpæn · ˌdʒæpəˈniːz}{le Japon · japonais(e)}{2}
\lexmot{jeans}{dʒiːnz}{un jean}{10}
\lexmot{juice}{dʒuːs}{le jus (de fruit)}{7}
\lexlettre{K}
\lexmot{a key}{ə ˈkiː}{une clé}{5}
\lexmot{a kitchen}{ə ˈkɪtʃɪn}{une cuisine}{8}
\lexlettre{L}
\lexmot{a lamp}{ə ˈlæmp}{une lampe}{5}
\lexmot{a language}{ə ˈlæŋɡwɪdʒ}{une langue}{2}
\lexmot{a laptop}{ə ˈlæptɑːp}{un ordinateur portable}{5}
\lexmot{last name}{ˈlæst ˌneɪm}{nom de famille}{3}
\lexmot{last night}{ˌlæst ˈnaɪt}{hier soir, la nuit dernière}{17}
\lexmot{last week}{ˌlæst ˈwiːk}{la semaine dernière}{17}
\lexmot{last weekend}{ˌlæst ˈwiːkend}{le week-end dernier}{18}
\lexmot{last year}{ˌlæst ˈjɪr}{l'année dernière}{17}
\lexmot{a latte}{ə ˈlɑːteɪ}{un café au lait}{11}
\lexmot{a laundromat}{ə ˈlɑːndrəmæt}{une laverie automatique}{8}
\lexmot{Léa speaking.}{ˈleɪə ˌspiːkɪŋ}{Léa à l'appareil.}{13}
\lexmot{leave}{liːv}{partir}{13}
\lexmot{leave · left}{liːv · left}{partir, quitter}{19}
\lexmot{Let's celebrate!}{ˌlets ˈseləbreɪt}{Fêtons ça !}{20}
\lexmot{Let's go out!}{ˌlets ɡoʊ ˈaʊt}{Sortons !}{15}
\lexmot{Let's meet at six.}{ˌlets ˈmiːt ət ˈsɪks}{Retrouvons-nous à six heures.}{15}
\lexmot{like}{laɪk}{aimer, apprécier}{7}
\lexmot{listening to music}{ˈlɪsənɪŋ tə ˈmjuːzɪk}{écouter de la musique}{7}
\lexmot{a little}{ə ˈlɪtəl}{un peu}{2}
\lexmot{a living room}{ə ˈlɪvɪŋ ˌruːm}{un salon, une salle de séjour}{8}
\lexmot{a long time ago}{ə ˌlɑːŋ ˈtaɪm əˈɡoʊ}{il y a longtemps}{17}
\lexmot{Look!}{lʊk}{Regarde ! / Regardez !}{4}
\lexmot{look for}{ˈlʊk fɚ}{chercher}{13}
\lexmot{a lot of}{ə ˈlɑːt əv}{beaucoup de}{10}
\lexmot{love}{lʌv}{adorer, aimer beaucoup}{7}
\lexmot{love · loved}{lʌv · lʌvd}{adorer}{18}
\lexmot{lovely}{ˈlʌvli}{charmant, ravissant}{8}
\lexlettre{M}
\lexmot{make · made}{meɪk · meɪd}{faire, fabriquer}{19}
\lexmot{make a wish}{ˌmeɪk ə ˈwɪʃ}{faire un vœu}{20}
\lexmot{maple syrup}{ˌmeɪpəl ˈsɪrəp}{le sirop d'érable}{18}
\lexmot{married}{ˈmerid}{marié(e)}{4}
\lexmot{Maybe another time.}{ˌmeɪbi əˌnʌðɚ ˈtaɪm}{Une autre fois, peut-être.}{15}
\lexmot{meat}{miːt}{la viande}{7}
\lexmot{meet}{miːt}{retrouver ; faire la connaissance de}{15}
\lexmot{meet · met}{miːt · met}{rencontrer}{19}
\lexmot{the menu}{ðə ˈmenjuː}{la carte, le menu}{11}
\lexmot{Mexico · Mexican}{ˈmeksɪkoʊ · ˈmeksɪkən}{le Mexique · mexicain(e)}{2}
\lexmot{milk}{mɪlk}{le lait}{7}
\lexmot{a mirror}{ə ˈmɪrɚ}{un miroir}{8}
\lexmot{mom · dad}{mɑːm · dæd}{maman · papa \registre{américain}}{4}
\lexmot{Monday}{ˈmʌndeɪ}{lundi}{6}
\lexmot{money}{ˈmʌni}{l'argent}{10}
\lexmot{mother · father}{ˈmʌðɚ · ˈfɑːðɚ}{la mère · le père}{4}
\lexmot{the mountains}{ðə ˈmaʊntənz}{la montagne}{16}
\lexmot{a movie}{ə ˈmuːvi}{un film}{15}
\lexmot{Mr.\ Carter}{ˈmɪstɚ ˈkɑːrtɚ}{monsieur Carter}{1}
\lexmot{Mrs.\ Lopez}{ˈmɪsɪz ˈloʊpez}{madame Lopez (femme mariée)}{1}
\lexmot{Ms.\ Clarke}{mɪz ˈklɑːrk}{madame Clarke (sans préciser)}{1}
\lexmot{a muffin}{ə ˈmʌfɪn}{un muffin}{11}
\lexmot{a museum}{ə mjuːˈziːəm}{un musée}{9}
\lexmot{a musical}{ə ˈmjuːzɪkəl}{une comédie musicale}{15}
\lexmot{My name is Jack.}{maɪ ˈneɪm ɪz ˈdʒæk}{Je m'appelle Jack.}{1}
\lexmot{my parents}{maɪ ˈperənts}{mes parents}{2}
\lexlettre{N}
\lexmot{nervous}{ˈnɝːvəs}{stressé, nerveux}{20}
\lexmot{new}{nuː}{neuf, nouveau}{5}
\lexmot{next month}{ˌnekst ˈmʌnθ}{le mois prochain}{16}
\lexmot{next week}{ˌnekst ˈwiːk}{la semaine prochaine}{14}
\lexmot{next year}{ˌnekst ˈjɪr}{l'année prochaine}{16}
\lexmot{Nice to meet you.}{ˌnaɪs tə ˈmiːt juː}{Enchanté(e).}{1}
\lexmot{Nice to meet you too.}{ˌnaɪs tə ˌmiːt juː ˈtuː}{Enchanté(e) moi aussi.}{1}
\lexmot{niece · nephew}{niːs · ˈnefjuː}{la nièce · le neveu}{4}
\lexmot{No, I can't.}{ˌnoʊ aɪ ˈkænt}{Non. (réponse courte)}{12}
\lexmot{No, I'm not.}{ˌnoʊ aɪm ˈnɑːt}{Non. (réponse courte)}{2}
\lexmot{noon · midnight}{nuːn · ˈmɪdnaɪt}{midi · minuit}{6}
\lexmot{not at all}{ˌnɑːt ət ˈɑːl}{pas du tout}{12}
\lexmot{Not bad.}{ˌnɑːt ˈbæd}{Pas mal.}{1}
\lexmot{not very well}{ˌnɑːt ˌveri ˈwel}{pas très bien}{12}
\lexmot{a notebook}{ə ˈnoʊtbʊk}{un carnet, un cahier}{5}
\lexmot{now}{naʊ}{maintenant}{13}
\lexlettre{O}
\lexmot{of course}{əv ˈkɔːrs}{bien sûr}{7}
\lexmot{Of course I can!}{əv ˌkɔːrs aɪ ˈkæn}{Bien sûr que je peux !}{12}
\lexmot{old}{oʊld}{vieux, ancien}{5}
\lexmot{on sale}{ɑːn ˈseɪl}{en promotion, en solde}{10}
\lexmot{on Saturday morning}{ɑːn ˌsætɚdeɪ ˈmɔːrnɪŋ}{samedi matin}{18}
\lexmot{on Sundays}{ɑːn ˈsʌndeɪz}{le dimanche (tous les dimanches)}{6}
\lexmot{on the corner}{ɑːn ðə ˈkɔːrnɚ}{au coin de la rue}{8}
\lexmot{on the corner of}{ɑːn ðə ˈkɔːrnɚ əv}{à l'angle de}{9}
\lexmot{on the left · on the right}{ɑːn ðə ˈleft · ɑːn ðə ˈraɪt}{sur la gauche · sur la droite}{9}
\lexmot{on weekends}{ɑːn ˈwiːkendz}{le week-end}{7}
\lexmot{one day}{ˌwʌn ˈdeɪ}{un jour}{16}
\lexmot{One more time, please.}{ˌwʌn mɔːr ˈtaɪm pliːz}{Encore une fois, s'il vous plaît.}{3}
\lexmot{a one-way ticket}{ə ˌwʌn weɪ ˈtɪkɪt}{un aller simple}{19}
\lexmot{an only child}{ən ˌoʊnli ˈtʃaɪld}{un enfant unique}{4}
\lexmot{orange}{ˈɔːrɪndʒ}{orange \otNuance{F0903B}}{5}
\lexmot{orange juice}{ˈɔːrɪndʒ ˌdʒuːs}{du jus d'orange}{11}
\lexlettre{P}
\lexmot{pack}{pæk}{faire ses bagages}{16}
\lexmot{pancakes}{ˈpænkeɪks}{des pancakes}{18}
\lexmot{pants}{pænts}{un pantalon \registre{américain}}{10}
\lexmot{parents}{ˈperənts}{les parents (père et mère)}{4}
\lexmot{a park}{ə ˈpɑːrk}{un parc}{9}
\lexmot{a party}{ə ˈpɑːrti}{une fête, une soirée}{15}
\lexmot{pasta}{ˈpɑːstə}{les pâtes}{7}
\lexmot{a pen}{ə ˈpen}{un stylo}{5}
\lexmot{a pencil}{ə ˈpensəl}{un crayon}{5}
\lexmot{a pharmacy}{ə ˈfɑːrməsi}{une pharmacie}{9}
\lexmot{a phone}{ə ˈfoʊn}{un téléphone}{5}
\lexmot{a phone number}{ə ˈfoʊn ˌnʌmbɚ}{un numéro de téléphone}{3}
\lexmot{a photo}{ə ˈfoʊtoʊ}{une photo}{4}
\lexmot{pink}{pɪŋk}{rose \otNuance{F29BC2}}{5}
\lexmot{a plan}{ə ˈplæn}{un projet}{16}
\lexmot{a plant}{ə ˈplænt}{une plante}{5}
\lexmot{play basketball}{ˌpleɪ ˈbæskɪtbɑːl}{jouer au basket}{12}
\lexmot{play soccer}{ˌpleɪ ˈsɑːkɚ}{jouer au football \registre{américain}}{12}
\lexmot{play tennis}{ˌpleɪ ˈtenɪs}{jouer au tennis}{12}
\lexmot{play the guitar}{ˌpleɪ ðə ɡɪˈtɑːr}{jouer de la guitare}{6}
\lexmot{play the piano}{ˌpleɪ ðə piˈænoʊ}{jouer du piano}{12}
\lexmot{playing sports}{ˈpleɪɪŋ ˈspɔːrts}{faire du sport}{7}
\lexmot{please}{pliːz}{s'il te plaît, s'il vous plaît}{1}
\lexmot{poor}{pʊr}{pauvre}{17}
\lexmot{a post office}{ə ˈpoʊst ˌɑːfɪs}{une poste}{9}
\lexmot{a present}{ə ˈprezənt}{un cadeau}{5}
\lexmot{proud}{praʊd}{fier}{20}
\lexmot{purple}{ˈpɝːpəl}{violet \otNuance{7B4FB3}}{5}
\lexlettre{Q}
\lexmot{quite well}{ˌkwaɪt ˈwel}{plutôt bien}{12}
\lexlettre{R}
\lexmot{reading}{ˈriːdɪŋ}{la lecture, lire}{7}
\lexmot{Really?}{ˈriːəli}{Vraiment ?}{2}
\lexmot{red}{red}{rouge \otNuance{E0574C}}{5}
\lexmot{relax}{rɪˈlæks}{se détendre}{18}
\lexmot{rent a car}{ˌrent ə ˈkɑːr}{louer une voiture}{16}
\lexmot{a reservation}{ə ˌrezɚˈveɪʃən}{une réservation}{19}
\lexmot{a restaurant}{ə ˈrestərɑːnt}{un restaurant}{9}
\lexmot{rice}{raɪs}{le riz}{7}
\lexmot{rich}{rɪtʃ}{riche}{17}
\lexmot{ride a bike}{ˌraɪd ə ˈbaɪk}{faire du vélo}{12}
\lexmot{right now}{ˌraɪt ˈnaʊ}{en ce moment même}{13}
\lexmot{a round-trip ticket}{ə ˌraʊnd trɪp ˈtɪkɪt}{un billet aller-retour}{19}
\lexmot{a rug}{ə ˈrʌɡ}{un tapis}{8}
\lexmot{run}{rʌn}{courir}{12}
\lexlettre{S}
\lexmot{a salad}{ə ˈsæləd}{une salade}{11}
\lexmot{a sandwich}{ə ˈsænwɪtʃ}{un sandwich}{7}
\lexmot{Saturday}{ˈsætɚdeɪ}{samedi}{6}
\lexmot{say · said}{seɪ · sed}{dire}{19}
\lexmot{a scarf}{ə ˈskɑːrf}{une écharpe}{10}
\lexmot{see · saw}{siː · sɑː}{voir}{19}
\lexmot{see you later}{ˌsiː juː ˈleɪtɚ}{à plus tard}{1}
\lexmot{see you tomorrow}{ˌsiː juː təˈmɑːroʊ}{à demain}{1}
\lexmot{a series}{ə ˈsɪriːz}{une série (télé)}{18}
\lexmot{a shelf}{ə ˈʃelf}{une étagère}{8}
\lexmot{She's on another call.}{ʃiz ɑːn əˌnʌðɚ ˈkɑːl}{Elle est en ligne.}{13}
\lexmot{a shirt}{ə ˈʃɝːt}{une chemise}{10}
\lexmot{shoes}{ʃuːz}{des chaussures}{10}
\lexmot{a show}{ə ˈʃoʊ}{un spectacle}{15}
\lexmot{sick}{sɪk}{malade}{12}
\lexmot{sing}{sɪŋ}{chanter}{12}
\lexmot{single}{ˈsɪŋɡəl}{célibataire}{4}
\lexmot{sister · brother}{ˈsɪstɚ · ˈbrʌðɚ}{la sœur · le frère}{4}
\lexmot{sit}{sɪt}{être assis, s'asseoir}{13}
\lexmot{six months ago}{ˌsɪks ˈmʌnθs əˈɡoʊ}{il y a six mois}{20}
\lexmot{ski}{skiː}{skier}{12}
\lexmot{a skirt}{ə ˈskɝːt}{une jupe}{10}
\lexmot{sleep}{sliːp}{dormir}{6}
\lexmot{small}{smɑːl}{petit}{5}
\lexmot{sneakers}{ˈsniːkɚz}{des baskets \registre{américain}}{13}
\lexmot{a sofa}{ə ˈsoʊfə}{un canapé}{8}
\lexmot{soon}{suːn}{bientôt}{16}
\lexmot{Sorry?}{ˈsɑːri}{Pardon ?}{3}
\lexmot{soup}{suːp}{de la soupe}{11}
\lexmot{a souvenir}{ə ˌsuːvəˈnɪr}{un souvenir (objet)}{19}
\lexmot{Spain · Spanish}{speɪn · ˈspænɪʃ}{l'Espagne · espagnol(e)}{2}
\lexmot{speak English}{ˌspiːk ˈɪŋɡlɪʃ}{parler anglais}{12}
\lexmot{spring}{sprɪŋ}{le printemps}{14}
\lexmot{the stairs}{ðə ˈsterz}{l'escalier}{8}
\lexmot{start · started}{stɑːrt · ˈstɑːrtɪd}{commencer}{18}
\lexmot{start work}{ˌstɑːrt ˈwɝːk}{commencer le travail}{6}
\lexmot{a station}{ə ˈsteɪʃən}{une gare}{9}
\lexmot{a statue}{ə ˈstætʃuː}{une statue}{17}
\lexmot{stay}{steɪ}{rester, séjourner}{16}
\lexmot{stay · stayed}{steɪ · steɪd}{rester}{18}
\lexmot{stay home}{ˌsteɪ ˈhoʊm}{rester à la maison}{18}
\lexmot{a steak}{ə ˈsteɪk}{un steak}{11}
\lexmot{stop · stopped}{stɑːp · stɑːpt}{(s')arrêter}{18}
\lexmot{a store}{ə ˈstɔːr}{un magasin}{10}
\lexmot{a stove}{ə ˈstoʊv}{une cuisinière}{8}
\lexmot{a street}{ə ˈstriːt}{une rue}{9}
\lexmot{study}{ˈstʌdi}{étudier}{6}
\lexmot{study · studied}{ˈstʌdi · ˈstʌdid}{étudier}{18}
\lexmot{a subway station}{ə ˈsʌbweɪ ˌsteɪʃən}{une station de métro}{9}
\lexmot{a suit}{ə ˈsuːt}{un costume}{13}
\lexmot{a suitcase}{ə ˈsuːtkeɪs}{une valise}{16}
\lexmot{summer}{ˈsʌmɚ}{l'été}{14}
\lexmot{Sunday}{ˈsʌndeɪ}{dimanche}{6}
\lexmot{a supermarket}{ə ˈsuːpɚˌmɑːrkɪt}{un supermarché}{9}
\lexmot{Sure!}{ʃʊr}{Bien sûr ! Volontiers !}{15}
\lexmot{a surprise party}{ə sɚˈpraɪz ˌpɑːrti}{une fête surprise}{20}
\lexmot{a sweater}{ə ˈswetɚ}{un pull}{10}
\lexmot{swim}{swɪm}{nager}{12}
\lexlettre{T}
\lexmot{a table}{ə ˈteɪbəl}{une table}{8}
\lexmot{take · took}{teɪk · tʊk}{prendre}{19}
\lexmot{take a shower}{ˌteɪk ə ˈʃaʊɚ}{prendre une douche}{6}
\lexmot{take photos}{ˌteɪk ˈfoʊtoʊz}{prendre des photos}{19}
\lexmot{take the subway}{ˌteɪk ðə ˈsʌbweɪ}{prendre le métro}{6}
\lexmot{talk to}{ˈtɑːk tə}{parler à}{13}
\lexmot{tall}{tɑːl}{haut, grand}{17}
\lexmot{a tea}{ə ˈtiː}{un thé}{11}
\lexmot{tea}{tiː}{le thé}{7}
\lexmot{a team}{ə ˈtiːm}{une équipe}{12}
\lexmot{a text}{ə ˈtekst}{un texto, un SMS}{13}
\lexmot{thank you}{ˈθæŋk juː}{merci}{1}
\lexmot{Thank you for everything!}{ˈθæŋk juː fɚ ˈevriθɪŋ}{Merci pour tout !}{20}
\lexmot{Thank you very much.}{ˈθæŋk juː ˌveri ˈmʌtʃ}{Merci beaucoup.}{1}
\lexmot{thanks}{θæŋks}{merci \registre{courant}}{1}
\lexmot{That sounds great!}{ðæt ˌsaʊndz ˈɡreɪt}{Ça a l'air super !}{15}
\lexmot{That's Jack.}{ˌðæts ˈdʒæk}{C'est Jack.}{3}
\lexmot{That's right.}{ˌðæts ˈraɪt}{C'est exact.}{2}
\lexmot{a theater}{ə ˈθiːətɚ}{un théâtre}{15}
\lexmot{think · thought}{θɪŋk · θɑːt}{penser}{19}
\lexmot{This is Emma.}{ˈðɪs ɪz ˈemə}{Voici Emma. / Je te présente Emma.}{1}
\lexmot{this weekend}{ðɪs ˈwiːkend}{ce week-end}{16}
\lexmot{Thursday}{ˈθɝːzdeɪ}{jeudi}{6}
\lexmot{a ticket}{ə ˈtɪkɪt}{un billet, une place}{15}
\lexmot{a tie}{ə ˈtaɪ}{une cravate}{13}
\lexmot{Time flies!}{ˌtaɪm ˈflaɪz}{Le temps passe vite !}{20}
\lexmot{a tip}{ə ˈtɪp}{un pourboire}{11}
\lexmot{To go, please.}{tə ˈɡoʊ pliːz}{À emporter, s'il vous plaît.}{11}
\lexmot{today}{təˈdeɪ}{aujourd'hui}{13}
\lexmot{tomorrow}{təˈmɑːroʊ}{demain}{16}
\lexmot{tonight}{təˈnaɪt}{ce soir}{15}
\lexmot{a tour}{ə ˈtʊr}{une visite guidée}{19}
\lexmot{a train station}{ə ˈtreɪn ˌsteɪʃən}{une gare}{19}
\lexmot{travel}{ˈtrævəl}{voyager}{16}
\lexmot{traveling}{ˈtrævəlɪŋ}{les voyages, voyager}{7}
\lexmot{a trip}{ə ˈtrɪp}{un voyage}{16}
\lexmot{try · tried}{traɪ · traɪd}{essayer, goûter}{18}
\lexmot{a T-shirt}{ə ˈtiːʃɝːt}{un tee-shirt}{10}
\lexmot{Tuesday}{ˈtuːzdeɪ}{mardi}{6}
\lexmot{turn left · turn right}{ˌtɝːn ˈleft · ˌtɝːn ˈraɪt}{tourner à gauche · à droite}{9}
\lexmot{a TV}{ə ˌtiːˈviː}{une télé}{5}
\lexmot{two days ago}{ˌtuː ˈdeɪz əˈɡoʊ}{il y a deux jours}{17}
\lexlettre{U}
\lexmot{the UK · British}{ðə ˌjuːˈkeɪ · ˈbrɪtɪʃ}{le Royaume-Uni · britannique}{2}
\lexmot{an umbrella}{ən ʌmˈbrelə}{un parapluie}{5}
\lexmot{the US · American}{ðə ˌjuːˈes · əˈmerɪkən}{les États-Unis · américain(e)}{2}
\lexlettre{V}
\lexmot{a vacation}{ə veɪˈkeɪʃən}{des vacances \registre{américain}}{16}
\lexmot{vegetables}{ˈvedʒtəbəlz}{les légumes}{7}
\lexmot{very well}{ˌveri ˈwel}{très bien}{12}
\lexmot{visit}{ˈvɪzɪt}{rendre visite à ; visiter}{16}
\lexmot{visit · visited}{ˈvɪzɪt · ˈvɪzɪtɪd}{visiter}{18}
\lexmot{voicemail}{ˈvɔɪsmeɪl}{la messagerie vocale}{13}
\lexlettre{W}
\lexmot{wait for}{ˈweɪt fɚ}{attendre}{13}
\lexmot{wake up}{ˌweɪk ˈʌp}{se réveiller}{6}
\lexmot{walk}{wɑːk}{marcher, aller à pied}{6}
\lexmot{walk · walked}{wɑːk · wɑːkt}{marcher, se promener}{18}
\lexmot{walking}{ˈwɑːkɪŋ}{la marche, se promener}{7}
\lexmot{want · wanted}{wɑːnt · ˈwɑːntɪd}{vouloir}{18}
\lexmot{a washing machine}{ə ˈwɑːʃɪŋ məˌʃiːn}{une machine à laver}{8}
\lexmot{a watch}{ə ˈwɑːtʃ}{une montre}{5}
\lexmot{watch · watched}{wɑːtʃ · wɑːtʃt}{regarder}{18}
\lexmot{watch TV}{ˌwɑːtʃ ˌtiːˈviː}{regarder la télé}{6}
\lexmot{watching movies}{ˈwɑːtʃɪŋ ˈmuːviz}{regarder des films}{7}
\lexmot{water}{ˈwɑːtɚ}{l'eau}{7}
\lexmot{wear}{wer}{porter (un vêtement)}{13}
\lexmot{Wednesday}{ˈwenzdeɪ}{mercredi}{6}
\lexmot{Welcome!}{ˈwelkəm}{Bienvenue !}{1}
\lexmot{wet}{wet}{mouillé, trempé}{14}
\lexmot{What a nice photo!}{ˌwʌt ə ˌnaɪs ˈfoʊtoʊ}{Quelle jolie photo !}{4}
\lexmot{What are you going to do?}{ˌwʌt ɑːr juː ˌɡoʊɪŋ tə ˈduː}{Qu'est-ce que tu vas faire ?}{16}
\lexmot{What can I get you?}{ˌwʌt kən aɪ ˈɡet juː}{Qu'est-ce que je vous sers ?}{11}
\lexmot{What color is it?}{ˌwʌt ˈkʌlɚ ɪz ɪt}{De quelle couleur est-il / elle ?}{5}
\lexmot{What did you do?}{ˌwʌt dɪd juː ˈduː}{Qu'est-ce que tu as fait ?}{18}
\lexmot{What kind of music?}{ˌwʌt ˈkaɪnd əv ˈmjuːzɪk}{Quel genre de musique ?}{7}
\lexmot{What size?}{ˌwʌt ˈsaɪz}{Quelle taille ?}{10}
\lexmot{What time is it?}{ˌwʌt ˈtaɪm ɪz ɪt}{Quelle heure est-il ?}{6}
\lexmot{What's the weather like?}{ˌwʌts ðə ˈweðɚ laɪk}{Quel temps fait-il ?}{14}
\lexmot{What's your email address?}{ˌwʌts jʊr ˈiːmeɪl ˌædres}{Quelle est ton adresse e-mail ?}{3}
\lexmot{What's your name?}{ˌwʌts jʊr ˈneɪm}{Comment tu t'appelles ? / Comment vous appelez-vous ?}{1}
\lexmot{What's your phone number?}{ˌwʌts jʊr ˈfoʊn ˌnʌmbɚ}{Quel est ton numéro de téléphone ?}{3}
\lexmot{When's your birthday?}{ˌwenz jʊr ˈbɝːθdeɪ}{C'est quand, ton anniversaire ?}{14}
\lexmot{Where are you from?}{ˌwer ɑːr juː ˈfrʌm}{D'où viens-tu ? / D'où venez-vous ?}{2}
\lexmot{Where were you?}{ˌwer wɚ ˈjuː}{Où étais-tu ?}{17}
\lexmot{Where were you born?}{ˌwer wɚ juː ˈbɔːrn}{Où es-tu né(e) ?}{17}
\lexmot{Where's the bank?}{ˌwerz ðə ˈbæŋk}{Où est la banque ?}{9}
\lexmot{white}{waɪt}{blanc \otNuance{FFFFFF}}{5}
\lexmot{Who's calling, please?}{ˌhuːz ˈkɑːlɪŋ ˌpliːz}{C'est de la part de qui ?}{13}
\lexmot{Who's that?}{ˌhuːz ˈðæt}{Qui est-ce ?}{3}
\lexmot{Who's this?}{ˌhuːz ˈðɪs}{Qui est-ce ? (sur une photo)}{4}
\lexmot{wife · husband}{waɪf · ˈhʌzbənd}{la femme (l'épouse) · le mari}{4}
\lexmot{a window}{ə ˈwɪndoʊ}{une fenêtre}{8}
\lexmot{winter}{ˈwɪntɚ}{l'hiver}{14}
\lexmot{Would you like something to drink?}{wʊd juː ˌlaɪk ˌsʌmθɪŋ tə ˈdrɪŋk}{Vous désirez quelque chose à boire ?}{11}
\lexmot{Would you like to come?}{wʊd juː ˌlaɪk tə ˈkʌm}{Tu veux venir ? Voulez-vous venir ?}{15}
\lexlettre{Y}
\lexmot{yellow}{ˈjeloʊ}{jaune \otNuance{F2C53D}}{5}
\lexmot{Yes, I am.}{ˌjes aɪ ˈæm}{Oui. (réponse courte)}{2}
\lexmot{Yes, I can.}{ˌjes aɪ ˈkæn}{Oui. (réponse courte)}{12}
\lexmot{Yes, please. · No, thank you.}{ˌjes ˈpliːz · ˌnoʊ ˈθæŋk juː}{Oui, volontiers. · Non, merci.}{11}
\lexmot{yesterday}{ˈjestɚdeɪ}{hier}{17}
\lexmot{yesterday morning}{ˌjestɚdeɪ ˈmɔːrnɪŋ}{hier matin}{17}
\lexmot{You too!}{ˌjuː ˈtuː}{Toi aussi ! / Vous aussi !}{1}
\lexmot{You're welcome.}{jʊr ˈwelkəm}{De rien. / Je vous en prie.}{1}
\end{lexique}
